\providecommand{\AnnotationOptions}{}
\expandafter\documentclass\expandafter[\AnnotationOptions]{lumenbook}
\title{Reading marks / 阅读标记}
\author{Lumenbook}
\begin{document}
\chapter{Reading with a pencil}
\section{Underlines and small annotations}
Underline a hypothesis: \textul{the boundary values vanish}.
中文也可以画线：\textul{先确认初始条件与边界条件，再判断结论的适用范围。}

\textwave{A condition to check / 尚待确认的条件} uses a wave.
Use \textstrike{an incorrect interpretation / 原来的错误理解} to record a revision.
\textaccent[LumenBlue]{Blue text / 蓝色文字} can distinguish a term.

The phrase \annotate{weak solution}{弱解首先指满足积分恒等式的函数。
这里的标记自动连接到编号边注。} links the underlined text to a numbered note.

\section{A palette for reading}
\textmark{Yellow / 黄色：关键结论}
\textmark[LumenMarkGreen]{Green / 绿色：已经理解}
\textmark[LumenMarkBlue]{Blue / 蓝色：概念与定义}
\textmark[LumenMarkRose]{Rose / 玫瑰色：待确认}
\textmark[LumenMarkViolet]{Violet / 紫色：联系与推广}

\textmark{HIGHLIGHT-START. This deliberately long bilingual passage wraps
over several lines. 阅读英文教材时，使用荧光底色标出结论，使用下划线标出条件，
在重新阅读时就能迅速区分论证的输入与输出。HIGHLIGHT-END.}

\begin{markblock}[LumenMarkGreen][An understood step / 已理解的步骤]
The boundary term vanishes because the boundary values are zero.
这一整段使用淡绿色背景，可以包含多段文字与数学内容。
\[
  E'(t)+\mathmark{\norm{\nabla u(t)}_{L^2}^2}=0.
\]
\end{markblock}
\begin{markblock}[LumenMarkRose][A question to revisit / 待回看的问题]
Which regularity assumption justifies integration by parts?
分部积分需要什么正则性条件？
\end{markblock}

\newpage
\section{Long passages and the margin}
\mn{\textul{边注中的下划线}。\par
  \textmark[LumenMarkGreen]{边注也可以用背景色标注中文或 English terms。}}
\textul{UNDERLINE-START. Mark the assumptions before reading the proof.
中文下划线也应当随文字自动换行，而不把整个句子锁成一个长盒子。
这一点对于正文和窄边栏都很重要。UNDERLINE-END.}

\begin{markblock}[LumenMarkViolet][Connections / 知识联系]
\begin{itemize}
  \item \textbf{Hypothesis.} Check the domain and boundary conditions.
  \item \textbf{Conclusion.} Separate existence, uniqueness and regularity.
  \item \textbf{Mechanism.} Reconstruct the energy argument.
\end{itemize}
\end{markblock}
\begin{codeblock}[language=LaTeX,title={A compact interface}]
\textul{中文与 English 下划线}
\textwave{有疑问的文字}
\textstrike{删除的文字}
\textmark[LumenMarkYellow]{黄色背景}
\annotate{weak solution}{这里写编号边注}
\begin{markblock}[LumenMarkBlue][标题可省略]
整段内容，支持列表、公式与跨页。
\end{markblock}
\end{codeblock}
\section{\textmark[LumenMarkYellow]{Marked headings}}
The PDF outline keeps this heading as ordinary searchable text.
\emph{Emphasis remains italic.} Native emphasis is not converted to an underline.
\end{document}
